\documentclass{article}
% Source: C:/Users/rober/Downloads/Topic_01_Teacher_Edition.pdf, PDF pages 91-92 (printed pages 57A-57B): Topic Assessment Form A with the TE's marked answers, Form B with marked answers, Item Analysis for both forms. Cross-checked against docs/a2-lesson-targets.md lines 282-310 (Topic 1 assessment audit, algebra2-live-worksheet repo).
\usepackage{amsmath}
\usepackage{amssymb}
\newenvironment{assessment-meta}{}{}   % topic, title, forms, item count, printed pages, sources
\newenvironment{item-analysis}{}{}     % Savvas's Item / DOK / Skills Review and Practice table (same for both forms)
\newenvironment{ta-item}[2]{}{}        % #1 = item number, #2 = form (A or B)
\newcommand{\answer}[1]{}              % the TE's marked answer, inside a ta-item
\newcommand{\placeholder}[2]{}         % #1 = type (graph, table, figure, number-line), #2 = complete description
\newcommand{\uncertain}[2]{#1}         % #1 = best reading, #2 = alternatives and why
\newcommand{\te}[1]{}                  % transcriber's note (a discrepancy, a publisher misprint, a computed check)
\begin{document}
\begin{assessment-meta}
topic: 1
title: Linear Functions and Systems
forms: A and B (parallel item for item)
items: 14 per form
printed: 57A-57B
sources: Topic_01_Teacher_Edition.pdf, PDF pages 91-92 (printed 57A Form A, 57B Form B and Item Analysis); cross-check: algebra2-live-worksheet/docs/a2-lesson-targets.md lines 282-310
\end{assessment-meta}
\te{TE page 57A callout, verbatim: Topic Assessment Assess students' understanding of topic concepts and skills using the Topic Assessment found at SavvasRealize.com. These auto-scored online assessments provide students with a breadth of technology-enhanced item types. There are two versions of the Topic Assessment, Form A and Form B. These two versions, available in print and at SavvasRealize.com, are parallel tests that assess the same content item for item. The Item Analysis chart on the next page can be used for both versions.}
\te{Running heads and footers agree: both TE pages head "TOPIC 1 Linear Functions and Systems", footers "TOPIC 1 | 57A | TOPIC ASSESSMENT" and "TOPIC 1 | 57B | TOPIC ASSESSMENT"; student pages head "1 Topic Assessment Form A" / "1 Topic Assessment Form B", footers "enVision Algebra 2 | 1 of 2 / 2 of 2 | Assessment Sourcebook". No running-head misprint on these pages.}

% Source page 92: 57B Item Analysis
\begin{item-analysis}
Item & DOK & Skills Review and Practice \\
1 & 2 & F25 \\
2 & 2 & F25 \\
3 & 2 & F31 \\
4 & 2 & F25, F56 \\
5 & 2 & F58 \\
6 & 2 & F22 \\
7 & 3 & F78, F79 \\
8 & 2 & F78, F79 \\
9 & 2 & F78, F79 \\
10 & 2 & F78, F79 \\
11 & 2 & A56 \\
12 & 2 & A12, A14 \\
13 & 2 & A59 \\
14 & 3 & A15 \\
\end{item-analysis}
\te{The printed table has no DOK 1 item: twelve items are DOK 2 and items 7 and 14 are DOK 3. The audit in docs/a2-lesson-targets.md (lines 288-292) lists items 1, 4, 5, 12 as DOK 1 and says there is no DOK 3 item; that DOK table disagrees with the TE. The TE table rules.}

% Source page 91: 57A Topic Assessment Form A (student page 1 of 2)
\begin{ta-item}{1}{A}
Use this graph for Items 1--3.
The graph represents Quinten's distance from the start line as he runs a cross country race. The x-axis represents the time, in hours, and the y-axis represents distance, in miles.
\placeholder{graph}{Shared graph for Items 1-3 (part of the prompt, black). Axes labeled x (right) and y (up) with arrowheads, origin O. x labels 2, 4, 6, 8 with grid lines every 1 unit; y labels 4, 8, 12, 16 with grid lines every 2 units; the grid runs about one cell left of the y-axis and below the x-axis, to x = 9 and y = 18. Piecewise-linear graph: closed dot at (0,0), segment to closed dot (2,4), segment to (5,7), horizontal segment from (5,7) to (6,7), segment to closed dot (8,13). No arrows on the graph.}
1. What is the range of the function?
[ [ ] , [ ] ]
\answer{1A. [0, 13] (boxes: 0, 13)}
\end{ta-item}

\begin{ta-item}{2}{A}
2. What does the domain of the function tell you about the race?
A. The race is 13 miles long.
B. Quinten takes 8 hours to finish the race.
C. Quinten takes a break at the 7-mile mark.
D. Quinten runs fastest in the last 2 hours.
\answer{2A. B}
\end{ta-item}

\begin{ta-item}{3}{A}
3. What is the average rate of change over the interval [0, 2]?
[ ] mi/h
\answer{3A. 2 mi/h}
\end{ta-item}

\begin{ta-item}{4}{A}
4. Over what interval is the graph of (y=-|x-2|) increasing?
A. ((2,\infty))
B. ((-\infty,2))
C. ((-\infty,-2))
D. ((-2,\infty))
\answer{4A. B}
\te{Choices print in two columns: A and B in the left column, C and D in the right.}
\end{ta-item}

\begin{ta-item}{5}{A}
5. Identify the translations of the parent function (f(x)=x^2) that give (g(x)=(x-5)^2-4).
A. up 4 units, left 5 units
B. up 4 units, right 5 units
C. down 4 units, left 5 units
D. down 4 units, right 5 units
\answer{5A. D}
\end{ta-item}

\begin{ta-item}{6}{A}
6. Graph the function.
(y=\begin{cases}3, \text{ if } -2\le x<-1\\ 2-x, \text{ if } -1<x<0\\ 2, \text{ if } 0\le x\le 4\end{cases})
\placeholder{graph}{Blank grid in the prompt: axes labeled x and y with arrowheads at both ends, origin O; x labels -4, -2, 2, 4 and y labels -4, -2, 2, 4; grid lines every 1 unit, window -5 to 5 on both axes. The magenta drawn answer: closed dot at (-2,3), horizontal segment to an open circle at (-1,3); segment from (-1,3) down to (0,2); horizontal segment at y = 2 from (0,2) to a closed dot at (4,2). No dot is drawn at (0,2).}
\answer{6A. Drawn answer: closed (-2,3) to open (-1,3); segment from open (-1,3) to (0,2); y = 2 from (0,2) to closed (4,2).}
\end{ta-item}

\begin{ta-item}{7}{A}
7. Enzo's profit over the last few months has been falling. In December he made \$900, in January his profit was \$500, and in March, he lost \$300. Write the recursive formula for the arithmetic sequence that represents Enzo's profit.
(a_1=) [ ]
(a_n=a_{n-1}+) [ ]
\answer{7A. (a_1=900), (a_n=a_{n-1}+(-400)) (boxes: 900, -400)}
\te{computed: 900, 500, 100, -300 for Dec, Jan, Feb, Mar; d = -400 agrees.}
\end{ta-item}

% Source page 91: 57A Topic Assessment Form A (student page 2 of 2)
\begin{ta-item}{8}{A}
8. The following recursive formula represents the number of prizes offered in each round of a game show.
(a_n=\begin{cases}2, \text{ if } n=1\\ a_{n-1}+7, \text{ if } n>1\end{cases})
Write the explicit formula for the number of prizes and use it to calculate the number of prizes in round 7.
(a_n=) [ ] + [ ] ((n-1))
[ ] prizes
\answer{8A. (a_n=2+7(n-1)); 44 prizes (boxes: 2, 7, 44)}
\end{ta-item}

\begin{ta-item}{9}{A}
9. Mya is testing out her new drone. She flies it 11 m and lands, then flies 22 m, then 33 m, and so on. The last test flight is 77 m. What is the total distance that Mya's drone travels during the tests?
A. 275 m
B. 308 m
C. 143 m
D. 385 m
\answer{9A. B}
\te{Choices print in two columns: A and B left, C and D right. computed: 11+22+...+77 = 7(11+77)/2 = 308; agrees.}
\end{ta-item}

\begin{ta-item}{10}{A}
10. In a theater, there are 12 chairs in the first row, each row has 2 more chairs than the previous row, and there are 10 rows. If (f(n)) represents the explicit formula for the number of seats in each row, what is a reasonable domain for (f)?
A. All positive integers
B. Multiples of 2, 12 through 30
C. Integers 1 through 10
D. Integers 0 through 10
\answer{10A. C}
\end{ta-item}

\begin{ta-item}{11}{A}
11. Use the graph to solve (|x+3|-1=(x+2)^2).
\placeholder{graph}{Part of the prompt (black). Axes labeled x and y with arrowheads at both ends, origin O; x labels -4, -2, 2, 4 and y labels 4, 2 (partly covered by the two curves where they cross near (0, 2)), -2, -4; grid lines every 1 unit, window -5 to 5 on both axes. Two curves, no points marked: the V graph of y = |x+3| - 1 with vertex (-3,-1), through (-4,0), (-2,0) and (0,2), arrows up-left past (-5,1) and up-right toward (3,5); the parabola y = (x+2)^2 with vertex (-2,0), through (-1,1) and (0,4), arrows at both upper ends. The graphs meet at (-2,0) and (-1,1).}
(x=) [ ], (x=) [ ]
\answer{11A. (x=-2), (x=-1) (boxes: -2, -1)}
\te{computed: x = -2 gives 0 = 0 and x = -1 gives 1 = 1; agrees.}
\end{ta-item}

\begin{ta-item}{12}{A}
12. Solve the system of equations.
(\begin{cases}x+4y=15\\ 2x-y=-15\end{cases})
(x=) [ ], (y=) [ ]
\answer{12A. (x=-5), (y=5) (boxes: -5, 5)}
\end{ta-item}

\begin{ta-item}{13}{A}
13. Solve the system of equations.
(\begin{cases}3x+2y+z=16\\ 4x-y=-5\\ y+z=11\end{cases})
A. ((0,5,6))
B. ((0,\frac{5}{3},2))
C. ((16,-5,11))
D. ((-2,11,6))
\answer{13A. A}
\te{Choices print in two columns: A and B left, C and D right.}
\end{ta-item}

\begin{ta-item}{14}{A}
14. Dwayne buys 1 loaf of bread and 6 eggs for \$4.45. Sheila buys 1 loaf of bread and 1 dozen eggs for \$5.65. Write and solve a system of equations to find how much the store charges for bread and for eggs. Then calculate how much Mina pays for 3 loaves of bread and 2 dozen eggs.
Bread costs \$ [ ] per loaf.
Eggs cost \$ [ ] per egg.
Mina pays \$ [ ].
\answer{14A. Bread costs \$3.25 per loaf. Eggs cost \$0.20 per egg. Mina pays \$14.55. (boxes: 3.25, 0.20, 14.55)}
\te{computed: b + 6e = 4.45, b + 12e = 5.65 give e = 0.20, b = 3.25; 3(3.25) + 24(0.20) = 14.55; agrees.}
\end{ta-item}

% Source page 92: 57B Topic Assessment Form B (student page 1 of 2)
\begin{ta-item}{1}{B}
Use this graph for Items 1--3.
The graph shows the distance of a remote control drone above the ground as it flies west to east. The x-axis represents the distance from a central point and the y-axis represents the distance above the ground, in m.
\placeholder{graph}{Shared graph for Items 1-3 (part of the prompt, black). Axes labeled x and y with arrowheads at both ends, origin O; x labels -4, -2, 2, 4 and y labels 2, 4; grid lines every 1 unit, window about -5 to 5 on x and -1 to 5 on y. Piecewise-linear graph: closed dot at (-4,0), segment up to (-2,2), segment down to (0,0), segment up to closed dot (4,3). No arrows on the graph.}
1. What is the range of the function?
[ [ ] , [ ] ]
\answer{1B. [0, 3] (boxes: 0, 3)}
\end{ta-item}

\begin{ta-item}{2}{B}
2. What does the domain of the function tell you about the drone's flight?
A. The drone landed once along its flight.
B. The drone flew up to 3 m high.
C. The drone traveled 8 m horizontally.
D. The drone flew for 8 minutes.
\answer{2B. C}
\end{ta-item}

\begin{ta-item}{3}{B}
3. What is the average rate of change over the interval [-4, 4]?
[ ] m/min
\answer{3B. 0.375 m/min}
\te{computed: (3 - 0)/(4 - (-4)) = 0.375; agrees. The printed unit label m/min does not fit the stem: the x-axis is horizontal distance in m, so the rate is metres of height per metre of horizontal distance. Transcribed as printed.}
\end{ta-item}

\begin{ta-item}{4}{B}
4. Over what interval is the graph of (y=-2+|x-2|) increasing?
A. ((2,\infty))
B. ((-\infty,2))
C. ((-\infty,-2))
D. ((-2,\infty))
\answer{4B. A}
\te{Choices print in two columns: A and B left, C and D right.}
\end{ta-item}

\begin{ta-item}{5}{B}
5. Identify the translations of the parent function (f(x)=|x|) that give (g(x)=2+|x+3|).
A. down 3 units, left 2 units
B. down 3 units, right 2 units
C. up 2 units, left 3 units
D. up 2 units, right 3 units
\answer{5B. C}
\end{ta-item}

\begin{ta-item}{6}{B}
6. Graph the function.
(y=\begin{cases}-2, \text{ if } 0\le x<1\\ 3-x, \text{ if } 1\le x\le 2\\ 1, \text{ if } 2<x\le 4\end{cases})
\placeholder{graph}{Blank grid in the prompt: axes labeled x and y with arrowheads at both ends, origin O; x labels -4, -2, 2, 4 and y labels -4, -2, 2, 4; grid lines every 1 unit, window -5 to 5 on both axes. The magenta drawn answer: closed dot at (0,-2), horizontal segment to an open circle at (1,-2); closed dot at (1,2), segment down to (2,1); horizontal segment at y = 1 from (2,1) to a closed dot at (4,1). No dot is drawn at (2,1).}
\answer{6B. Drawn answer: closed (0,-2) to open (1,-2); closed (1,2) to (2,1); y = 1 from (2,1) to closed (4,1).}
\end{ta-item}

\begin{ta-item}{7}{B}
7. The wait time at the DMV increases for each new person who arrives. Jane arrived first and waited 1 minute. Pete arrived next and waited 1.2 minutes. Kal arrives fourth and waits 1.6 minutes. Write the recursive formula for the arithmetic sequence that represents the wait time.
(a_1=) [ ]
(a_n=a_{n-1}+) [ ]
\answer{7B. (a_1=1), (a_n=a_{n-1}+0.2) (boxes: 1, 0.2)}
\te{computed: 1, 1.2, 1.4, 1.6; d = 0.2 agrees.}
\end{ta-item}

% Source page 92: 57B Topic Assessment Form B (student page 2 of 2)
\begin{ta-item}{8}{B}
8. The following recursive formula represents the number of players left after week (n) of a reality show.
(a_n=\begin{cases}15, \text{ if } n=1\\ a_{n-1}-2, \text{ if } n>1\end{cases})
Write the explicit formula for the number of players left and use it to calculate how many players are left after 6 weeks.
(a_n=) [ ] + [ ] ((n-1))
[ ] players
\answer{8B. (a_n=15+(-2)(n-1)); 5 players (boxes: 15, -2, 5)}
\te{computed: a_6 = 15 - 2(5) = 5; agrees.}
\end{ta-item}

\begin{ta-item}{9}{B}
9. Manny's Pizza shows declining sales over the last few months. Their monthly net profit in hundreds of dollars follows the arithmetic sequence 4, 1, -2, -5, \ldots, -56. What is their total profit over this time period?
A. -\$77,700
B. -\$54,600
C. -\$54,200
D. -\$49,000
\answer{9B. B}
\te{computed: 21 terms, sum 21(4 + (-56))/2 = -546 hundred = -\$54,600; agrees.}
\end{ta-item}

\begin{ta-item}{10}{B}
10. In a theater, there are 20 chairs in the first row, each row has 4 more chairs than the previous row, and there are 15 rows. If (f(n)) represents the explicit formula for the number of seats in each row, what is a reasonable domain for (f)?
A. All positive integers
B. Multiples of 4, 20--76
C. Integers 1--15
D. Integers 0--15
\answer{10B. C}
\end{ta-item}

\begin{ta-item}{11}{B}
11. Use the graph to solve (\frac{1}{2}|x+1|+1=|x|).
\placeholder{graph}{Part of the prompt (black). Axes labeled x and y with arrowheads at both ends, origin O; x labels -4, -2, 2, 4 and y labels -4, -2, 2, 4; grid lines every 1 unit, window -5 to 5 on both axes. Two V graphs, no points marked: y = |x| with vertex at the origin, arrows up-left past (-5,5) and up-right past (5,5); y = (1/2)|x+1| + 1 with vertex (-1,1), through (3,3), arrows up-left toward (-5,3) and up-right toward (5,4). The graphs meet at (-1,1) and (3,3).}
(x=) [ ], (x=) [ ]
\answer{11B. (x=-1), (x=-3) (boxes: -1, -3)}
\te{computed: the printed key x = -3 is wrong. At x = -3 the left side is 2 and |x| = 3. The graph's crossings and the algebra give x = -1 and x = 3. Same finding as docs/a2-lesson-targets.md line 303. Printed key kept.}
\end{ta-item}

\begin{ta-item}{12}{B}
12. Solve the system of equations.
(\begin{cases}2x-y=7\\ x-y=10\end{cases})
(x=) [ ], (y=) [ ]
\answer{12B. (x=-3), (y=-13) (boxes: -3, -13)}
\end{ta-item}

\begin{ta-item}{13}{B}
13. Solve the system of equations.
(\begin{cases}x+3y+2z=1\\ x-z=4\\ -4x+3y=-12\end{cases})
A. ((6,-1,-1))
B. ((29,5,-11))
C. ((3,0,-1))
D. ((1,4,-12))
\answer{13B. C}
\te{Choices print in two columns: A and B left, C and D right. computed: (3,0,-1) satisfies all three; agrees.}
\end{ta-item}

\begin{ta-item}{14}{B}
14. Emery makes 8 necklaces and 3 bracelets, using a total of 172 beads. Then she makes 1 more necklace and 5 more bracelets, using 77 beads. Write and solve a system of equations to find how many beads Emery uses in each necklace and bracelet. Then calculate how many beads Emery uses to make 4 necklaces and 7 bracelets.
Emery uses [ ] beads for each necklace and [ ] beads for each bracelet.
Emery uses [ ] beads.
\answer{14B. 17 beads for each necklace and 12 beads for each bracelet; Emery uses 152 beads. (boxes: 17, 12, 152)}
\te{computed: 8n + 3b = 172, n + 5b = 77 give n = 17, b = 12; 4(17) + 7(12) = 152; agrees.}
\end{ta-item}

\end{document}
